\begin{frame}
    \frametitle{Discussion and Policy Implications}
    \begin{itemize}[itemsep=0.65em]
      \item Informal discretion affects the regulatory environment alongside written zoning rules.
      \item Local representation can provide neighborhood information, while local decision-making authority may undervalue housing benefits elsewhere in the city.
      \item This paper abstracts from the aggregate supply effects but shows individual local actors matter 
    \end{itemize}
\end{frame}

\begin{frame}
    \frametitle{Conclusion}
    \textbf{Summary of findings}
    \begin{itemize}[itemsep=0.25em,topsep=0.2em,parsep=0pt,partopsep=0pt]
      \item Aldermen create substantial differences in the regulatory environment within the same city.
      \item More stringent aldermen $\Rightarrow$ lower permitting and lower multifamily density on the more-stringent side of ward boundaries.
      \item These supply constraints appear to feed into rents.
    \end{itemize}

    \vspace{0.4em}
    \textbf{Implications}
    \begin{itemize}[itemsep=0.25em,topsep=0.2em,parsep=0pt,partopsep=0pt]
      \item Informal regulations can also affect housing supply, not just formal rules on the books.
      \item Reforms intended to increase housing supply should also target decision-making authority to be effective.
    \end{itemize}
\end{frame}
